\documentclass[11pt]{article}
\usepackage[T1]{fontenc}
\usepackage{lmodern,amsmath,amssymb,amsthm,booktabs,microtype}
\usepackage[margin=26mm]{geometry}
\usepackage[hidelinks]{hyperref}
\newtheorem{theorem}{Theorem}
\newtheorem{proposition}[theorem]{Proposition}
\newtheorem{lemma}[theorem]{Lemma}
\newtheorem{corollary}[theorem]{Corollary}
\theoremstyle{remark}\newtheorem{remark}[theorem]{Remark}
\newcommand{\Q}{\mathbb Q}\newcommand{\Z}{\mathbb Z}\newcommand{\C}{\mathbb C}
\newcommand{\F}{\mathbb F}\newcommand{\Gal}{\operatorname{Gal}}
\newcommand{\Spec}{\operatorname{Spec}}\newcommand{\tr}{\operatorname{tr}}
\newcommand{\workbenchurl}{https://github.com/KokunoYumeto/zeta-and-l-functions-workbench}
\newcommand{\recordurl}{https://github.com/KokunoYumeto/mathematics-workbenches}
\title{Cubic Artin representations and the scope of Borger's descent programme}
\author{}\date{24 November 2026}
\begin{document}\maketitle
\begin{abstract}
Borger proved that rational cohomology of schemes with commuting Frobenius lifts is abelian and becomes Tate over a cyclotomic field. We identify the tensor category generated by the Tate-normalized cohomology of smooth proper such schemes, and exhibit the obstruction to the cubic representation associated with the cube root of 269. We distinguish this obstruction from the existence of its finite arithmetic fiber and of relative Frobenius lifts over the Eisenstein field. A finite-order phase also fails to separate certain exterior-power summands even over rational coefficients. Finally, finite-ring averages identify the arithmetic cancellation still required by the cubic Riemann hypothesis. The arguments and text were produced by GPT-6 Astra in ChatGPT Pro mode; the phase calculation was independently re-proved by a separate instance of the same model, and no human has checked the text.
\end{abstract}
\section{Main statements and scope}
This note distinguishes the existence of a representation from an estimate for its Frobenius traces. Put $K=\Q(\omega)$, where $\omega^2+\omega+1=0$, and put $F=\Q(\alpha)$ and $E=K(\alpha)$, where $\alpha^3=269$. The group $\Gal(E/\Q)$ is $S_3$. Write $\rho$ for its two-dimensional sum-zero permutation representation. Its Artin function $L(s)=L(s,\rho)=\zeta_F(s)/\zeta(s)$ is entire; its Riemann hypothesis is not proved here.

\begin{theorem}\label{th:category}
Fix a prime $\ell$. In the category of finite-dimensional $\Q_\ell$-representations of $G_\Q$, the tensor category generated by
\[
 H^{2j}_{\mathrm{et}}(X_{\overline\Q},\Q_\ell)(j)
\]
for integers $j\ge0$ and smooth proper Borger $\Lambda$-schemes $X$ over rings $\Z[1/N]$ with an integer $N\ge1$, allowing direct sums, duals, and subquotients, is exactly the category of representations with finite abelian image. All odd cohomology groups of these $X$ vanish. The representation $\rho$ does not belong to this category.
\end{theorem}
Here $N$ and $X$ may vary; only finitely many primes are inverted, as in Borger (2009, Theorem 0.1). The converse is a statement about the generated category, not a claim that every representation is itself the whole cohomology of a single connected scheme.


\subsection*{Arithmetic fibers and relative Frobenius lifts}
Set $R=\Z[1/(3\cdot269)]$ and $A=R[X]/(X^3-269)$.
The degree-three fiber is $\Spec A$; the splitting cover is the
normalization of $\Spec R$ in $E$. The cyclic cubic cover is the
normalization of $\Spec\mathcal O_K[1/(3\cdot269)]$ in $E$.
A good prime ideal is one not dividing $3\cdot269$.
\begin{proposition}\label{th:frobenius}
At an unramified rational prime $p$ of $F$, the fixed degree-three fiber admits a rational Frobenius lift if and only if $p$ splits completely in $F$; that lift is unique. The fixed splitting cover over $R$ has the same property. Over $\mathcal O_K[1/(3\cdot269)]$, the cyclic cubic cover admits a unique relative Frobenius lift at every good prime ideal, and all such lifts commute.
\end{proposition}

\subsection*{Cubic phases}
For an integer $n\ge1$, let $V=\bigoplus_{j=0}^{2n}V_j$ over
$\Q(\omega)$, with every $V_j\ne0$. Let $U$ act on $V_j$ by
$\omega^{2j-2n}$, and let $e_W$ be the projector onto
$V_0\oplus V_{2n}$. The middle summands mean
$\bigoplus_{j=1}^{2n-1}V_j$.
\begin{proposition}\label{th:phase}
For $n\ge2$, no polynomial in $U$ and $U^{-1}$ over $\Q$ or $\Q(\omega)$ equals $e_W$. For $n=1$, one has
\[
 e_W=\frac{2-U-U^2}{3},
\]
and the least denominator at 3 in $\Q[U]$ is exactly 3. If equality is required only after quotienting by the middle summands, the identity already works and no denominator is necessary.
\end{proposition}

\begin{lemma}\label{lem:commuting-phase}
Let $V$ be a finite-dimensional complex vector space and $U,F$ commuting endomorphisms with $U^3=1$. Suppose $V^{U=1}=0$. If $B>0$, $C>0$, and
\[
 |\tr(UF^a)|\le C B^a\qquad(a\ge1),
\]
then every eigenvalue of $F$ has absolute value at most $B$. More generally, with all three eigenspaces allowed, the same conclusion holds provided no generalized eigenspace of $F$ has equal multiplicities for the three eigenvalues of $U$.
\end{lemma}

\subsection*{Finite-ring kernel averages}
For a positive integer $M$, put $A_{F,M}=\mathcal O_F/M\mathcal O_F$
and define
\[
 J_F(M)=\frac1{|A_{F,M}|}\sum_{a\in A_{F,M}}
        \log|\ker(m_a:A_{F,M}\longrightarrow A_{F,M})|.
\]
Here $m_a$ denotes multiplication by $a$; define $J_\Q$ in the
same way. If $p\mathcal O_F=\prod_i\mathfrak p_i^{e_i}$ with
residue degrees $f_i$, then for every integer $k\ge1$,
\[
 \frac{J_F(p^k)}{\log p}=\sum_i f_i\sum_{r=1}^{ke_i}p^{-f_i r}.
\]
At an unramified cubic prime the values of $J_F(p)/\log p$ are
$3/p$, $1/p+2/p^2$, and $3/p^3$ for the splitting types
$1,(12),(123)$, respectively. For
$M_y=\prod_{y<p\le2y}p$, $H=\lceil y^{3/2}\rceil$, and $y>269$,
\[
 \frac1H\sum_{a=A+1}^{A+H}\log\gcd(|a^3-269|,M_y)
   =J_F(M_y)+O(y^{-1/2}\log(2y)),
\]
with an effective constant independent of the integer $A$.
The complete Euler-factor differences obtained by removing one good
prime from $-L'(1)/L(1)$ are, in the same three splitting types,
\[
 \log p\left\{\frac2{p-1},\frac2{p^2-1},
                     -\frac{p+2}{p^2+p+1}\right\}.
\]
Their dyadic sum differs from $J_F(M_y)-J_\Q(M_y)$ by
$O(\log y/y)$. These identities do not establish a lower bound
$J_F(M_y)-J_\Q(M_y)\ge-C_\varepsilon y^{-1/2+\varepsilon}$
for every $\varepsilon>0$ and all sufficiently large $y$.
The required prime-distribution estimate is separate from the
unconditional inequality $J_F(M_y)\ge0$.

\subsection*{Kummer factors}
The fields and characters in the next two statements are local to
those statements; in particular their radicand need not be 269.
\begin{proposition}\label{prop:prime-kummer}
Let $r$ be an odd prime and $p$ a rational prime. Put $F=\Q(p^{1/r})$, $K_r=\Q(\zeta_r)$ and $E_r=K_r(p^{1/r})$. Let $\chi$ be any faithful character of $N=\Gal(E_r/K_r)$. Then, with all finite Euler factors included,
\[
 \frac{\zeta_F(s)}{\zeta(s)}=L_{K_r}(s,\chi),\qquad
 \zeta_{E_r}(s)=\zeta_{K_r}(s)L_{K_r}(s,\chi)^{r-1}.
\]
Under RH for Dirichlet $L$-functions, RH for either $\zeta_F$ or $\zeta_{E_r}$ is equivalent to RH for this one Hecke function.
\end{proposition}

\begin{proposition}\label{prop:composite-kummer}
Let $n\ge1$, $a\in\Q^\times$, $K_n=\Q(\zeta_n)$ and $E_n=K_n(a^{1/n})$. Put $m=[E_n:K_n]$, so that $m\mid n$, and choose a faithful character $\chi$ of the cyclic group $N=\Gal(E_n/K_n)$. Then
\[
 \zeta_{E_n}(s)=\zeta_{K_n}(s)
 \prod_{\substack{d\mid m\\d>1}}
 L_{K_n}\bigl(s,\chi^{m/d}\bigr)^{\varphi(d)}.
\]
The product is empty if $m=1$. Under Dirichlet RH, RH for $\zeta_{E_n}$ is equivalent to RH for every displayed Hecke factor. No assumption that $m=n$ is needed.
\end{proposition}

\subsection*{A growing CM comparison}
Fix $0<\epsilon<1/52$, put $\vartheta=\epsilon/100$,
$L=x^\epsilon$ and $\delta=x^{-\epsilon/8}$, and let $x$ be
sufficiently large. Choose a prime
$q\equiv5\pmod {12}$ with
$2x^{2+8\epsilon}\le q\le4x^{2+8\epsilon}$, and set
\[
 H_x=\Q(\zeta_{3q}),\qquad n=(q-1)/2,\qquad f=2n.
\]
Choose $m$ to be the nearest positive even integer to
$x^{1/2-\epsilon/16}$, $k$ the nearest integer to
$x^{9/2+10\epsilon}$, and $g_0=g_1=\lceil(ke^L)^8\rceil$.
The field $H_x$ is CM, contains $\omega$, and has degree $4n$.
These choices satisfy
\begin{align*}
 x^{-\vartheta}\sqrt{x\delta}&\le m\le x^\vartheta\sqrt{x\delta},&
 \log f&\le\sqrt{m/2-1},\\
 (ke^L)^8&\le g_i,&\log g_i&\le x^\vartheta L,\\
 x^{2+8\epsilon-\vartheta}&\le n\le x^{2+8\epsilon+\vartheta},&
 x^{9/2+10\epsilon-\vartheta}&\le k\le x^{9/2+10\epsilon+\vartheta},
\end{align*}
as well as $2\le m\le f\le k$ and $m+3\le k$.
Only these numerical conditions are asserted; the conductor,
place-set, radius, interpolation, partner-count and chain-of-lifts
conditions required by a theta construction are not asserted.

Let $\chi$ be either nontrivial cubic character of $E/K$, and put
$\chi_x=\chi\circ N_{H_x/K}$. It is nontrivial and has
\[
 \begin{gathered}
 \mathfrak f(\chi_x)=(269)\mathcal O_{H_x},\qquad
 |D_{H_x}|=3^{q-1}q^{2q-4},\\
 Q(\chi_x)=3^{q-1}q^{2q-4}269^{2(q-1)}.
 \end{gathered}
\]
There is no unramified degree-one prime of $H_x$ over any rational
prime $p\le2x$ with $p\nmid3q$. In particular,
\[
 \tr\!\left(\operatorname{Ind}_{G_{H_x}}^{G_\Q}\chi_x
                      (\operatorname{Frob}_p)\right)=0
       \quad(p\le2x,\ p\nmid3q\cdot269).
\]
The factorization
\[
 L_{H_x}(s,\chi_x)
   =\prod_{\eta\in\widehat{\Gal(H_x/K)}}L_K(s,\chi\eta)
\]
contains the fixed cubic factor at $\eta=1$, but the other factors
cancel its first-prime trace in this range. This is an obstruction
for this particular family, not an estimate for the fixed cubic sum.

\subsection*{Cyclotomic cofactors and projected coboundaries}
\begin{proposition}\label{prop:cofactor-prime}
Let $r$ be an odd prime, $p$ a rational prime, and $\ell\ne p$ a prime with $\ell\equiv1\pmod r$. Write $d=\operatorname{ord}_\ell(p)$. Then
\[
 p^{(\ell-1)/r}\not\equiv1\pmod\ell
 \quad\Longleftrightarrow\quad
 v_r(d)=v_r(\ell-1).
\]
In particular, if $m\ge2$, $\ell\mid\Phi_m(p)$ and $\ell\nmid m$, then $\operatorname{ord}_\ell(p)=m$. Such a prime divisor is an auxiliary-prime witness precisely when $r\mid m$ and $v_r(\ell-1)=v_r(m)$.
\end{proposition}

\begin{proposition}\label{prop:cofactor-algebra}
Let $q$ be a prime power, $r$ an odd prime dividing $q-1$, and $z\in\mathbf F_q^\times$. Put
\[
 A=\mathbf F_q[X]/(X^r-z),\qquad
 \alpha=X\bmod(X^r-z),\qquad \tau(\alpha)=\zeta_r\alpha,
\]
where $\zeta_r\in\mathbf F_q$ has exact order $r$. For $c\in A^\times$, set $b(c)=\tau(c)c^{-1}$. Then
\[
 z\notin(\mathbf F_q^\times)^r
 \quad\Longleftrightarrow\quad
 \text{some }c\in A^\times\text{ has }b(c)^{q-1}\ne1.
\]
In the nonpower case $A$ is the field $\mathbf F_{q^r}$, and $c=\alpha+1$ is such a witness. Every projected coboundary $b(c)^{q-1}$ has order dividing $\Phi_r(q)/r$, which is prime to $r$ and to $q-1$. In contrast, the existence of a nonscalar unprojected coboundary of order prime to $r$ does not imply the nonpower condition.
\end{proposition}

\section*{Proof overview}
Borger's theorem makes each representation abelian and a sum of Tate characters after cyclotomic restriction. Smooth proper purity forces the Tate character in degree $2j$ to be exactly the $j$th one. Cyclotomic zero-dimensional schemes then supply every finite abelian representation as a summand of a regular representation. The commutator subgroup of $S_3$ acts nontrivially on $\rho$, so it cannot occur. None of this prevents a nonabelian finite etale cover from producing $\rho$: that cover simply lacks rational Frobenius lifts at nonsplit primes. Over $K$, its cyclic Galois group instead gives commuting relative lifts.


For the phase statements, equal eigenvalues prevent the desired
projector, whereas nonzero phase multiplicities make large Frobenius
eigenvalues visible as poles of a trace generating function. The
finite-ring identities follow from capped valuations and the Chinese
remainder theorem; counting polynomial roots in intervals gives the
uniform sampling error.

The Kummer factorizations follow by inducing characters of the cyclic
translation group and retaining the local determinants at ramified
primes. In the growing cyclotomic family, the large modulus excludes
degree-one primes in the target range; induction then forces the
displayed trace to vanish. Finally, cyclic finite-field groups give
the cofactor valuation test and the projected-coboundary criterion.
These arguments identify exact representations and coefficients;
none supplies the missing prime-distribution bound.

\section{The exact category generated by smooth proper descent}
This section proves Theorem~\ref{th:category} with its coefficient and smoothness hypotheses intact. Borger (2009, Theorem 0.1 and Corollary 6.2) gives abelian Galois action and, after restriction to $G_{\Q(\zeta_N)}$, a direct sum of representations $\Q_\ell(-m)$. In smooth proper degree $i$, Frobenius eigenvalues have weight $i$. At a good prime splitting in $\Q(\zeta_N)$, an eigenvalue from $\Q_\ell(-m)$ has absolute value $p^m$. Purity gives $2m=i$. Consequently odd cohomology is zero, and in degree $2j$ every summand has $m=j$. Alternatively odd vanishing follows from Borger's diagonal Hodge numbers and Hodge decomposition. Twisting by $(j)$ therefore makes the open subgroup $G_{\Q(\zeta_N)}$ act trivially. The resulting image is finite and abelian.

A finite collection of finite abelian representations factors through one finite abelian quotient. Direct sums, tensor products, duals and subquotients continue to factor through that quotient, so every object generated in the theorem has finite abelian image. Conversely, let $V$ have finite abelian image. By Kronecker--Weber its fixed field lies in some $\Q(\zeta_N)$. The scheme
\[
 X_N=\Spec\Z[\zeta_N,1/N]
\]
is finite etale, hence smooth and proper, over $\Z[1/N]$. For $p\nmid N$, its Frobenius lift sends $\zeta_N$ to $\zeta_N^p$; for $p\mid N$ choose the identity. These operations commute, and the Frobenius condition at an inverted prime is vacuous. This is Borger's example in Remark 6.8. Its geometric degree-zero cohomology is the regular representation of $(\Z/N\Z)^\times$. By Maschke's theorem every $\Q_\ell$-representation of this finite group is a summand of a finite sum of regular representations. This includes $V$ and proves the converse.

The commutator subgroup of $S_3$ is $A_3$. A three-cycle on $\rho$ has characteristic polynomial $X^2+X+1$, so it is not the identity over $\Q_\ell$, including $\ell=3$. Every finite abelian representation kills the commutator subgroup. This proves the final assertion. The rational-coefficient hypothesis matters: no assertion about semisimplicity over $\Z_\ell$ is made.

\section{Arithmetic fibers, rational quotients and relative lifts}
This section constructs the cubic representation while showing why it does not contradict the preceding theorem. Set $R=\Z[1/(3\cdot269)]$. The fiber
\[
 A=R[X]/(X^3-269)
\]
of the finite etale map $x\mapsto x^3$ on the multiplicative group is finite etale of rank three. The polynomial is Eisenstein at 269 and its discriminant is $-27\cdot269^2$, a nonsquare. Thus its splitting group is $S_3$. The direct image of the constant $\Q_\ell$-sheaf is the three-letter permutation representation $\mathbf1\oplus\rho$; projection by subtracting the average realizes $\rho$ as a lisse sheaf. This is a finite-etale and degree-zero cohomology construction, with no proposed correspondence left to define.

For completeness a rational coordinate gives the same fiber. With $x^3=a\ne0$ set $r=x-x^{-1}$. Then
\[
 r^3+3r=a-a^{-1},\qquad x=\frac{a-r}{r^2+1}.
\]
The denominator is nonzero in characteristic zero when $a=269$. These identities identify the cubic fields defined by the two equations. The polynomial $r^3+3r$ is the Dickson polynomial $D_3(r,-1)$. For the three conjugate $x_i$, the ratio of Vandermonde determinants is
\[
 \prod_{i<j}\frac{(x_i-x_i^{-1})-(x_j-x_j^{-1})}{x_i-x_j}
 =\prod_{i<j}\left(1+\frac1{x_ix_j}\right)=1+a^{-2}.
\]
The extra order index is supported on $a^2+1=2\cdot97\cdot373$. After inverting $2a$, its valuations at 97 and 373 are one. Hence the nonnormal Dickson order must not be used to compute Artin Euler factors at these two primes. At such a prime $a^2=-1$ and the Dickson cubic becomes $(r-a)^2(r+2a)$, whereas the normalized original cubic splits completely: $(-a)^3=a$ and $p\equiv1\pmod3$.

\begin{proof}[Proof of Proposition~\ref{th:frobenius}]
Every endomorphism of the generic field $F$ is the identity, since the other two roots of its defining cubic are nonreal. Its reduction is the $p$th-power map on the entire finite etale fiber precisely when all residue fields are $\F_p$. Torsion-freeness gives uniqueness on the integral algebra. On the splitting cover, a Frobenius lift would have to act as the Frobenius at every prime over $p$. These Frobenius elements form one conjugacy class. Since a good decomposition group acts faithfully on the residue field, they must all coincide. Their common value is central in $S_3$, whose center is trivial. Thus again complete splitting is necessary and sufficient.

For a prime ideal $\mathfrak p$ of $K$ of norm $q$ prime to $3\cdot269$, the group of third roots of unity embeds in its residue field. There is a unique $\zeta_{\mathfrak p}\in\mu_3$ with
\[
 \zeta_{\mathfrak p}\equiv269^{(q-1)/3}\pmod{\mathfrak p}.
\]
The automorphism $\alpha\mapsto\zeta_{\mathfrak p}\alpha$ fixes the base and reduces to the $q$th-power map, because $\alpha^{q-1}=269^{(q-1)/3}$. The cyclic deck group is abelian, proving commutation and uniqueness.
\end{proof}

The traces at a good rational prime are $2,0,-1$ for the cycle types $1,(12),(123)$, respectively. They give Euler inverse factors $(1-T)^2$, $1-T^2$, and $1+T+T^2$. These are precisely the traces needed by the signed prime race. Existence of the sheaf identifies them prime by prime; it does not bound their aggregate cancellation.

There is also a direct finite-orbit obstruction in the Dickson family. If $r=z-z^{-1}$ has finite orbit under $D_3$, then $r(z^{3^i})=r(z^{3^j})$ for some $j>i$. The two points over a value of $w-w^{-1}$ are $w$ and $-w^{-1}$. Thus either $z^{3^j-3^i}=1$ or $z^{3^j+3^i}=-1$. In both cases $z$ is a root of unity. Conversely a root of unity gives a finite orbit under all odd Dickson operations. A finite Galois-stable set of such parameters and their inverse cubic fibers therefore lies in a cyclotomic field and has abelian monodromy. The parameter $269-269^{-1}$ escapes this restriction by having infinite orbit.

\section{A correction to the proposed cubic-unit projector}
This section gives a precise spectral obstruction that denominators cannot remove. Let $V=\bigoplus_{j=0}^{2n}V_j$ over $\Q(\omega)$, with every $V_j\ne0$, and let $U$ act on $V_j$ by $\omega^{2j-2n}$. These are the eigenvalues of $[\omega]^*$ on the usual exterior-power decomposition in degree $2n$. Let $e_W$ be the projector onto $V_0\oplus V_{2n}$.

\begin{proof}[Proof of Proposition~\ref{th:phase}]
For $n\ge2$, the middle index $j=3$ is available whenever $2n\ge4$, and its eigenvalue equals that at $j=0$. The case $n=2$ has $j=3<4$ as required. Any polynomial in $U,U^{-1}$ acts by the same scalar on these two nonzero summands, whereas $e_W$ acts by 1 and 0. This is a contradiction over any coefficient field. The analogous collision at $j=2n-3$ applies to the other extreme. For $n=1$ the three eigenvalues are $\omega,1,\omega^2$; the displayed expression is 1 on the first and third and 0 on the middle. Since all three occur, the minimal polynomial is $T^3-1$. Every polynomial has a unique representative of degree at most two modulo this monic polynomial. The displayed coefficients $2/3,-1/3,-1/3$ establish the exact denominator. The final assertion follows because identity induces identity on the quotient by the middle summands.
\end{proof}

Thus a finite function $c(n)$ describing the denominator of the proposed exact unit projector does not exist for $n\ge2$. The obstruction precedes any 3-adic estimate. A nonunit endomorphism with distinct relevant eigenvalues is a different question.

\subsection{Commuting cubic phases and reflected traces}
The obstruction to an exact projector does not invalidate every cubic phase estimate. The distinction is spectral and can be stated completely.

\begin{proof}[Proof of Lemma~\ref{lem:commuting-phase}]
For an eigenvalue $\beta$ of $F$, let $m_j(\beta)$ be the dimension of its generalized eigenspace on which $U$ acts as $\omega^j$. Simultaneous triangularization on each $U$-eigenspace gives
\[
 \tr(UF^a)=\sum_\beta c_\beta\beta^a,
 \qquad c_\beta=m_0(\beta)+m_1(\beta)\omega+m_2(\beta)\omega^2.
\]
The relation $1+\omega+\omega^2=0$ is the only rational relation among these three numbers. Thus $c_\beta=0$ precisely when all three multiplicities are equal. If $m_0(\beta)=0$ and the generalized eigenspace is nonzero, then
$\Re c_\beta=-(m_1(\beta)+m_2(\beta))/2<0$.
The generating function of the traces is
\[
 \sum_{a\ge1}\tr(UF^a)z^a
 =\sum_{\beta\ne0}\frac{c_\beta\beta z}{1-\beta z}.
\]
The assumed bound makes the left side holomorphic for $|z|<1/B$. Distinct nonzero eigenvalues give distinct simple poles on the right, and their coefficients are nonzero under the stated hypothesis. Therefore $|\beta|\le B$.
\end{proof}

Both qualifications are necessary. On the regular representation of $C_3$, take $F=L\,\mathrm{id}$ with arbitrary $L>0$. Then $\tr(UF^a)=0$ for every $a$, because the three phase multiplicities are equal. For a reflected trace take
\[
 U=\operatorname{diag}(\omega,\omega^{-1}),\qquad
 S=\begin{pmatrix}0&1\\1&0\end{pmatrix},\qquad F=L\,\mathrm{id}.
\]
Here $SUS^{-1}=U^{-1}$ and $\tr(SF^a)=0$ for every $a$, again for arbitrary $L$. This example is already compatible with the standard finite $S_3$ representation at split Frobenius elements after a Tate twist. In contrast, bounds for $\tr((UF)^a)$ with $UF=FU$ always bound the eigenvalues of $F$, because eigenvalues of $UF$ have the same absolute values and their ordinary power sums have positive integer multiplicities.

Over $K$, the cubic character lines are separately stable, and their deck operator commutes with Frobenius. Over $\Q$, transposition Frobenius elements exchange the lines. In a genuine representation of a cubic metaplectic cover, its central $\mu_3$ acts by a scalar and commutes with every operator, so the lemma itself presents no spectral obstruction. However, pairing a theta coefficient with a conjugate theta coefficient cancels that central character. The fixed-radicand correlation is an operation on Fourier indices; it is not supplied by multiplying the theta representation by its central scalar. The remaining task is to realize the analytic remainder of the prime estimate as a trace with controlled multiplicities. A finite arithmetic fiber alone does not identify that remainder.

\section{What an integral kernel average still has to prove}
This section identifies the numerical deficit left by integral nonnegativity. Let $M_y=\prod_{y<p\le2y}p$ and $A_F=\mathcal O_F/M_y\mathcal O_F$. Define
\[
 J_F(M_y)=\frac1{|A_F|}\sum_{a\in A_F}\log|\ker(m_a:A_F\to A_F)|.
\]
Writing $p\mathcal O_F=\prod_i\mathfrak p_i^{e_i}$ with residue degrees $f_i$, the Chinese remainder theorem and the tail-sum formula give exactly
\[
 \frac{J_F(p^k)}{\log p}=\sum_i f_i\sum_{r=1}^{ke_i}p^{-f_i r}.
\]
Indeed the probability that a uniform element is divisible by $\mathfrak p_i^r$ is $p^{-f_i r}$, and each unit of capped valuation contributes $f_i\log p$ to the logarithmic kernel size. At a good cubic prime,
\[
 J_F(p)=\log p\begin{cases}3/p&\text{complete split},\\1/p+2/p^2&\text{type }(12),\\3/p^3&\text{type }(123).\end{cases}
\]
For $H=\lceil y^{3/2}\rceil$ and $y>269$, root counting in an interval gives
\[
 \frac1H\sum_{a=1}^{H}\log\gcd(|a^3-269|,M_y)
 =J_F(M_y)+O(y^{-1/2}\log(2y)).
\]
Each of at most three roots modulo $p$ occurs $H/p+O(1)$ times; summing the error $O(\log p/H)$ over $p\le2y$, and the residue-degree terms $O(\log p/p^2)$, proves the assertion. The constant is effective and independent of the interval's starting point.

For the cubic target the missing lower bound is
\[
 J_F(M_y)-J_\Q(M_y)\ge-C_\varepsilon y^{-1/2+\varepsilon}
\]
with the relevant all-large-$y$ or sparse-window quantifiers proved in the separate analytic criterion. Its complete Euler-factor version is the change of $-L'(1)/L(1)$ on removing the dyadic primes: the per-prime differences are
\[
 \log p\left\{\frac2{p-1},\frac2{p^2-1},-\frac{p+2}{p^2+p+1}\right\}.
\]
They follow by summing $f_i/(p^{f_i}-1)$ and subtracting $1/(p-1)$. The finite kernel and complete expressions differ by $O(\log y/y)$.

Integrality alone yields $|\ker m_a|\ge1$, hence $J_F(M_y)\ge0$. It therefore gives only a relative lower bound $-J_\Q(M_y)$, of constant order, whereas the displayed error tends to zero. The three local cases have the same field degree and very different kernel averages. An invariant stable tangent degree, without a further comparison involving their distribution, does not supply the missing cancellation estimate. This identifies a needed estimate; it is not a theorem that all possible determinant arguments must fail.

\section{The Hecke factors of Kummer towers}
This section identifies precisely which nonabelian factors remain after Dirichlet Riemann hypotheses are granted. The argument is Artin formalism; it does not require a conjectural automorphic construction.

\begin{proof}[Proof of Proposition~\ref{prop:prime-kummer}]
The valuation of $p$ at a prime of $K_r$ is either $1$ or $r-1$, and hence is not divisible by $r$. Thus $[E_r:K_r]=r$. The splitting group is $G=C_r\rtimes\F_r^\times$, acting on $\F_r$ by affine transformations. Let $P$ be its permutation representation and $W=P-\mathbf1$. Restricted to the translation subgroup $N$, the representation $W$ is the sum of all nontrivial characters of $N$. The multiplicative subgroup acts transitively on these characters. It follows directly, by inducing a character of $N$, that
\[
 W\simeq\operatorname{Ind}_N^G\chi.
\]
The regular representation satisfies
\[
 \mathbf C[G]\simeq\operatorname{Ind}_N^G\mathbf1
       \oplus (r-1)W.
\]
For example, its dimension is $(r-1)+(r-1)^2=r(r-1)$, and the equality follows also by inducing the decomposition of $\mathbf C[N]$ into its $r$ characters. The two zeta identities now follow from the permutation interpretation of zeta functions and compatibility of Artin $L$-functions with induction. The inducing subgroup for $W$ is the translation subgroup $N$, not the multiplicative stabilizer of a point.

These identities include ramified primes: on restricting an induced representation to a decomposition group, its inertia invariants decompose according to primes of the intermediate field. Frobenius acts on each coset orbit by a cyclic block. The determinant of $1-T$ times such a block is the corresponding local determinant with $T$ replaced by the residue-degree power. Their product is exactly the local Euler factor over the intermediate field. This proves the local form of induction as well.

Every nontrivial finite-order Hecke function in the identities is entire. In the open critical strip a zero of one factor cannot be canceled by a pole of another. Since $\zeta_{K_r}$ is a product of Dirichlet functions, the claimed equivalences follow.
\end{proof}

\begin{proof}[Proof of Proposition~\ref{prop:composite-kummer}]
The extension $E_n/\Q$ is Galois because it is the splitting field of $X^n-a$. The group $N$ embeds into $\mu_n$ by $\sigma\mapsto\sigma(a^{1/n})/a^{1/n}$, so it is cyclic of order $m\mid n$. For $G=\Gal(E_n/\Q)$, induction of the regular character of $N$ gives
\[
 \mathbf C[G]\simeq\bigoplus_{j=0}^{m-1}\operatorname{Ind}_N^G\chi^j.
\]
Consequently $\zeta_{E_n}=\prod_{j=0}^{m-1}L_{K_n}(s,\chi^j)$. The $j=0$ factor is $\zeta_{K_n}$. Conjugation by a lift of $u\in(\Z/n\Z)^\times$ acts on $N$ by raising to the $u$th power. Reduction from $(\Z/n\Z)^\times$ onto $(\Z/m\Z)^\times$ is surjective, as follows prime by prime from the Chinese remainder theorem. Thus the characters of $N$ of a fixed order $d$ form one conjugacy orbit. Their Hecke functions over $K_n$ agree, since conjugation permutes prime ideals while preserving their norms and the ramified local determinants. There are $\varphi(d)$ such characters. Grouping the preceding product proves the formula. The argument uses no splitting of the group extension $1\to N\to G\to\Gal(K_n/\Q)\to1$. Entireness again proves the RH equivalence.
\end{proof}

This is the family of Dedekind hypotheses appearing in the usual conditional approach to Artin's primitive-root problem. For $r=3$ the fixed value $p=269$ is one member of that family. A theorem for this one field gives no uniform assertion for varying $p$, for composite $n$, or for all non-pure cubic fields. Any complexity consequence that needs a small auxiliary prime uniformly in an input $p$ must therefore be stated for a uniform family of Hecke hypotheses, rather than for the single numerical example.

\section{A growing CM family and the loss of degree-one primes}
This section tests a concrete family against both numerical size and prime distribution. Its purpose is to separate a feasible choice of degrees from a construction of the required theta lifts.
Fix $0<\epsilon<1/52$, put $\vartheta=\epsilon/100$, $L=x^\epsilon$ and $\delta=x^{-\epsilon/8}$. For every sufficiently large $x$, choose a prime
\[
 q\equiv5\pmod {12},\qquad
 2x^{2+8\epsilon}\le q\le4x^{2+8\epsilon}.
\]
Such a prime exists by the prime number theorem in the fixed progression modulo 12. Define
\[
 H_x=\Q(\zeta_{3q}),\quad n=(q-1)/2,\quad f=2n,
\]
and choose $m$ to be the nearest positive even integer to $x^{1/2-\epsilon/16}$ and $k$ the nearest integer to $x^{9/2+10\epsilon}$. Put $g_0=g_1=\lceil(ke^L)^8\rceil$. Then $H_x$ is CM, contains $\omega$, and has degree $4n$. The five numerical constraints used in the chain-length comparison hold for all sufficiently large $x$:
\begin{align*}
 x^{-\vartheta}\sqrt{x\delta}&\le m\le x^{\vartheta}\sqrt{x\delta},&
 \log f&\le\sqrt{m/2-1},\\
 (ke^L)^8&\le g_i,&\log g_i&\le x^\vartheta L,\\
 x^{2+8\epsilon-\vartheta}&\le n\le x^{2+8\epsilon+\vartheta},&
 x^{9/2+10\epsilon-\vartheta}&\le k\le x^{9/2+10\epsilon+\vartheta}.
\end{align*}
Indeed, fixed multiplicative constants and rounding errors are absorbed by $x^{\vartheta}$; $\log f=O_\epsilon(\log x)$ whereas $\sqrt m$ is a positive power of $x$; and
$8(L+\log k)/(x^{\vartheta}L)\to0$.
The additional elementary rank requirements $2\le m\le f\le k$ and $m+3\le k$ follow from the same exponents. This verifies the numerical subcollection only. The other conditions on local conductors, the set of places, lattice radii, interpolation nodes, partner counts, and the chain of lifts are not asserted for this family.

Let $\chi_x=\chi\circ N_{H_x/K}$ be the base change of the cubic character. It remains nontrivial: $[H_x:K]=q-1$ is prime to 3, whereas $[E:K]=3$. At the primes over 269, the extension $H_x/K$ is unramified and the cubic character is tamely ramified; its conductor exponent remains one. At every other finite prime it is unramified. Thus
\[
 \mathfrak f(\chi_x)=(269)\mathcal O_{H_x},\qquad
 \log N_{H_x/\Q}\mathfrak f(\chi_x)=2(q-1)\log269.
\]
The discriminant and Artin conductor over $\Q$ are, respectively,
\[
 |D_{H_x}|=3^{q-1}q^{2q-4},\qquad
 Q(\chi_x)=3^{q-1}q^{2q-4}269^{2(q-1)}.
\]
The discriminant formula follows from the cyclotomic discriminant formula, or from the product of the primitive Dirichlet conductors. At bounded height, or at height bounded by any fixed power of $x$, the logarithmic analytic conductor grows on the scale $x^{2+8\epsilon}\log x$, although the root discriminant grows only polynomially in $x$.

There is a separate, exact prime-norm obstruction. For every rational prime $p\le2x$ with $p\nmid3q$, the residue degree in $H_x$ is the order of $p$ in $(\Z/3q\Z)^\times$. It cannot be one, because $q>2x$ and $p\equiv1\pmod q$ would force $p\ge q+1$. Consequently this family has no unramified degree-one primes over the rational primes in the target range. Equivalently,
\[
 \tr\!\left(\operatorname{Ind}_{G_{H_x}}^{G_\Q}\chi_x
                  (\operatorname{Frob}_p)\right)=0
 \quad(p\le2x,\ p\nmid3q\cdot269).
\]
To prove the trace identity directly, $G_{H_x}$ is normal in $G_\Q$ and an induced character vanishes outside it. Frobenius belongs to $G_{H_x}$ precisely at completely split primes.

Induction over $K$ also gives the exact factorization
\[
 L_{H_x}(s,\chi_x)=
 \prod_{\eta\in\widehat{\Gal(H_x/K)}}L_K(s,\chi\eta).
\]
The factor $\eta=1$ is the target function. Its first-prime coefficient is canceled by the other factors throughout the range $p\le2x$, by character orthogonality. At the exceptional prime 3 the relative extension $H_x/K$ is unramified, its Frobenius is $3\bmod q\ne1$, and the same orthogonality applies. At 269 every twisted character remains ramified, so each local factor is 1. Inert rational primes already have zero first-prime coefficient over $K$. Thus the unweighted prime-ideal trace of this natural growing family does not recover the fixed cubic prime sum. A successful chain must supply an explicit character projection or a different family with the required prime-norm comparison. This obstruction concerns this particular cyclotomic choice, not every growing CM family.

\section{What a cyclotomic cofactor witness certifies}
This section makes precise two elementary algebraic tests that can occur in an auxiliary-prime construction. It separates their exact field-theoretic content from any running-time claim for finding or factoring a cofactor.

\begin{proof}[Proof of Proposition~\ref{prop:cofactor-prime}]
The multiplicative group of $\mathbf F_\ell$ is cyclic. Its subgroup of $r$th powers has order $(\ell-1)/r$, so it contains $p$ exactly when $d$ divides that order. Since $d\mid\ell-1$, removing one factor $r$ prevents divisibility precisely when their $r$-adic valuations agree.

For the cyclotomic assertion, $X^m-1$ is squarefree modulo $\ell$ because $\ell\nmid m$. Its factors $\Phi_e(X)$, $e\mid m$, therefore have disjoint root sets. A root of $\Phi_m$ cannot have smaller order, since it would then also be a root of the product for a proper divisor. Thus its order is $m$, and the first assertion applies. Primes dividing $m$ require their own local order calculation; they are not included in this last shortcut.
\end{proof}

For example,
\[
 \Phi_3(269)=72631=13\cdot37\cdot151.
\]
All three factors see $269$ with order three. The primes $13$ and $151$ are witnesses, since their multiplicative group orders have exactly one factor of three. The prime $37$ is not: $v_3(36)=2$, so an element of order three is a cube there. Divisibility of a cyclotomic value alone does not distinguish these cases.

There is a separate issue with a proposed correspondence between a nontrivial Kummer Frobenius and a nonscalar coboundary. Its unprojected form is false.

\begin{proof}[Proof of Proposition~\ref{prop:cofactor-algebra}]
Since all $r$th roots of unity lie in $\mathbf F_q$, the algebra is either split, when $z$ is an $r$th power, or a field of degree $r$. Indeed Frobenius permutes the roots by multiplication by one $r$th root of unity; its orbit has size one or $r$, and a fixed root is equivalent to the power condition. In the split case $A\simeq\mathbf F_q^r$, so every unit, hence every coboundary, has $(q-1)$st power one.

In the field case, the kernel of $c\mapsto b(c)$ is $\mathbf F_q^\times$. Thus its image has order
\[
 \frac{q^r-1}{q-1}=\Phi_r(q)
\]
and is exactly the norm-one subgroup. The gcd of this order with $q-1$ is $r$, since $\Phi_r(q)\equiv r\pmod{q-1}$. Also $v_r(\Phi_r(q))=1$: writing $q=1+r^a u$, $a\ge1$, the binomial expansion gives $\Phi_r(q)\equiv r\pmod{r^2}$. Raising to $q-1$ therefore gives a cyclic image of order $\Phi_r(q)/r$, prime to both $r$ and $q-1$.

The element $\alpha+1$ is nonzero. If
$(\zeta_r\alpha+1)/(\alpha+1)$ were a scalar $v\in\mathbf F_q$, linear independence of $1,\alpha$ would give both $v=\zeta_r$ and $v=1$, a contradiction. Thus this coboundary is not fixed by the $q$-power map and its $(q-1)$st power is not one.

For the final counterexample take $q=7$, $r=3$, $z=1$, and $\zeta_3=2$. In the split algebra choose
\[
 c=X^2+2X+5,\qquad b=5X^2+6X+2.
\]
At the roots $1,2,4$ the values of $c$ are $1,-1,1$, so $c$ is a unit. Direct multiplication modulo $X^3-1$ gives $bc=\tau(c)$ and $b^2=1$, while $b$ is not scalar. This is a nonscalar coboundary of order two even though the Kummer algebra is completely split. Its projected value $b^6$ is one, as the corrected criterion requires.
\end{proof}

The projection is therefore essential in this precise formulation. In the field
\[
 \mathbf F_7[X]/(X^3-2),
\]
the projected coboundary from $X+1$ is $1+X+X^2$, of exact order 19. These finite witnesses establish the algebraic distinction. A particular cofactor algorithm must additionally specify its projection, its admissible factorization and its search bounds before its eligible input primes or complexity can be characterized.

\section{Source dictionary and open comparisons}
The preceding constructions concern ordinary schemes and their cohomology. SGA 4, Expose XII, Theorem 5.1 is proper base change for an arbitrary Cartesian square: for sheaves of sets it gives the degree-zero isomorphism; for sheaves of groups it gives degree-one injectivity, upgraded to bijectivity for ind-finite groups; for torsion abelian sheaves it gives isomorphisms in every degree. Corollary 5.2 specializes to a geometric point and identifies the stalk with cohomology of the fiber, with the same separate coefficient hypotheses. Neither statement asserts the integral K-theory comparison or the number-field RH estimate. The passage to $\ell$-adic coefficients additionally requires the usual inverse-limit formalism.

Over $\Z[\omega]$, the normalized-ideal obstruction for imaginary quadratic fields has no force: every ideal $J=(b)$ is already principal, and every complex embedding has $|b|^2=N(J)$. Thus $F=K$ itself realizes the normalized generator, without any binomial theorem. This checks the conclusion for the Eisenstein field directly. It does not provide the growing CM host family, the required chain lengths, or a cubic theta remainder identity.

The open arithmetic comparisons are: construct the commuting cubic phase on the precise global theta space and identify its effective trace; prove a uniform elliptic-unit or determinant estimate with the conductor and prime-norm losses required by the window; or prove the one-sided finite-kernel bound. The finite cover, its Euler data, the ray-theta period, and the generated abelian category each have a precise role, but none of their existence statements establishes that bound.

\section*{Method and records}
The arguments and text were produced by GPT-6 Astra in ChatGPT Pro mode. A separate instance of the same model independently re-proved Proposition~\ref{th:phase} from its statement. Source reading and further checks are recorded at \href{\recordurl}{the permanent records}; programs and the detailed working records are at \href{\workbenchurl}{the workbench}. No human has checked this text. The original text and programs are dedicated to the public domain under CC0 1.0 Universal.

\begin{thebibliography}{99}
\bibitem{borger} Borger, J. (2009). \emph{$\Lambda$-rings and the field with one element}. arXiv:0906.3146. Theorem 0.1, Corollary 6.2, Remark 6.8. \url{https://arxiv.org/abs/0906.3146}
\bibitem{bs} Borger, J., \& de Smit, B. (2008). Galois theory and integral models of $\Lambda$-rings. \emph{Bulletin of the London Mathematical Society, 40}, 439--446. \url{https://doi.org/10.1112/BLMS/BDN024}
\bibitem{sga4} Artin, M., Grothendieck, A., \& Verdier, J.-L. (Eds.). (1973). \emph{Theorie des topos et cohomologie etale des schemas}, tome 3. Lecture Notes in Mathematics 305. Springer. Expose XII, Theoreme 5.1 and Corollaire 5.2.
\bibitem{deligne} Deligne, P. (1974). La conjecture de Weil. I. \emph{Publications Mathematiques de l'IHES, 43}, 273--307. \url{https://doi.org/10.1007/BF02684373}
\bibitem{milne} Milne, J. S. (2020). \emph{Class field theory} (Version 4.03). Chapter VIII, Section 10, pp.~261--263, for Artin induction; Chapter VI, Section 4, for progression density. \url{https://www.jmilne.org/math/CourseNotes/CFTc.pdf}
\end{thebibliography}
\end{document}
